\section{Methods}
\label{sec:methods}

\subsection{Study design and data source}

We conducted a retrospective cohort study of ED encounters with documented pain assessments using the Medical Information Mart for Intensive Care IV (MIMIC-IV) Emergency Department module.
MIMIC-IV contains de-identified clinical data from a tertiary academic medical center; calendar years are reported as de-identified anchor years.
The study was designed to characterize factors associated with time to first pain reassessment across routine ED care rather than to test a single demographic hypothesis.
This report follows STROBE guidance for observational cohort studies \citep{vonElm2007STROBE}.

\subsection{Study population}

We identified ED stays with at least one documented initial pain score and eligible ED diagnoses.
The analytic cohort was restricted to encounters with ED diagnoses of acute pancreatitis or trauma (including fall, fracture/dislocation, and other trauma subtypes), reflecting a clinically homogeneous population with expected pain documentation and reassessment workflows.
Figure~\ref{fig:flow} summarizes cohort construction.

Starting from \FlowAllED{} ED stays in the extract, we applied diagnosis-based inclusion, excluded small or unknown race/ethnicity categories, required a positive initial pain score with valid time-to-event information, excluded undocumented insurance, and required a documented Emergency Severity Index (ESI) triage level.
The resulting primary analytic cohort comprised \CohortN{} encounters.

\subsection{Outcome and follow-up}

The primary outcome was time from the \textbf{initial documented pain score} to the \textbf{first subsequent pain reassessment}, measured in minutes.
Patients without a documented reassessment before ED departure were censored at the end of the ED stay.
A reassessment was defined as any pain score recorded strictly after the initial pain documentation timestamp.

For descriptive reporting, we also calculated the proportion of patients reassessed within 60 and 120 minutes.
In sensitivity analyses, we examined post-analgesic reassessment timing with time zero set to the first analgesic administration after initial pain documentation (Section~\ref{sec:sensitivity}).

\subsection{Covariates}

Covariates were organized into four sequential domains informed by clinical plausibility and care-pathway timing (conceptual diagram, Supplementary Figure~S1):

\begin{itemize}
  \item \textbf{M1 (clinical presentation):} initial pain score and injury/diagnosis group (acute pancreatitis reference).
  \item \textbf{M2 (demographic and social):} race/ethnicity (White reference), insurance (private reference), preferred language (English reference when available), age group, and sex.
  \item \textbf{M3 (illness severity):} ESI triage level, initial vital signs (heart rate, respiratory rate, systolic blood pressure, z-scored), and comorbidity count when ICD-based flags were available.
  \item \textbf{M4 (ED workflow context):} arrival mode (walk-in reference), arrival shift (day reference), weekend arrival, ED crowding proxies (arrivals in the prior hour and census at the hour of initial pain documentation), and policy era (5-year de-identified year bins).
\end{itemize}

Disposition (admitted vs.\ discharged home) and analgesic treatment were considered downstream of initial pain documentation and were \textbf{not} included in the primary sequential models.
Vital-sign coefficients were estimated in M3 but are reported in the supplement because they are adjusted for but not emphasized in the primary forest plot.

\subsection{Statistical analysis}

We used Cox proportional hazards models to estimate hazard ratios (HRs) for time to first reassessment.
HR $>1$ indicates faster reassessment; HR $<1$ indicates slower reassessment.
We fit four sequential models (M1--M4) adding covariate domains in the order above.
\textbf{M4} was prespecified as the primary adjusted model for inference on demographic, insurance, and workflow associations.

Additional analyses included:
\begin{enumerate}
  \item \textbf{Sequential attenuation plots} for race, insurance, and initial pain across M1--M4;
  \item \textbf{Within-acuity models} using an M4-style specification within ESI strata (1--2, 3, and 4--5);
  \item \textbf{Severe-pain sensitivity} restricting to initial pain scores 7--10 (and pain score 10);
  \item \textbf{Temporal trend analysis} comparing unadjusted and M4-adjusted probabilities of reassessment within 60 minutes across policy eras, with a supplementary continuous-year Cox model;
  \item \textbf{Post-analgesic pathway analysis} with time zero at first analgesic;
  \item \textbf{Disposition pathway sensitivity} comparing M4 with M4 plus admitted-vs-home disposition;
  \item \textbf{Inverse probability of treatment weighting (IPTW)} comparing Medicaid vs.\ private insurance, with propensity scores estimated using M1--M3 covariates only.
\end{enumerate}
\label{sec:sensitivity}

Analyses were conducted in Python 3 using \texttt{lifelines} and \texttt{statsmodels}; code to reproduce all tables and figures is available in the study repository (\texttt{scripts/run\_analysis.py}).

\subsection{Ethics}

MIMIC-IV is a publicly available, de-identified database; credentialed users complete required training and data use agreements.
Institutional review board approval was not required for secondary analysis of de-identified data per the data use agreement.
